% GENERATED FILE. Edit canonical lessons, then run npm run generate:books.
% canonical-source-hash: fnv1a64:f9ec692036725e80
% canonical-lessons: FA-C525-naqqash, FA-C525-bum, FA-C525-qalam-mu, FA-C525-abrang, FA-C525-galeri

\chapter{Painting}
\label{ch:fa-painting}
\hlchaptermodality{525}

\begin{chapteropening}
\textbf{By the end of this chapter:} \emph{I can talk about painting.}

Complete the last lesson of chapter 525: I can talk about painting.
\end{chapteropening}

\section[\texorpdfstring{\fa{نقاش} (naqqâsh)}{نقاش (naqqâsh)}]{\fa{نقاش} (naqqâsh) — a painter}

\label{lesson:FA-C525-naqqash}



\begin{quote}
Before the new one: say the Persian for to play a role, then the Persian for an actor.
\end{quote}

\subsection*{You'll want to know: \fa{نقاش}}
\textbf{\fa{نقاش}} (\emph{naqqâsh}) — \textquotedblleft{}a painter\textquotedblright{}.

\begin{quote}
\fa{این} \fa{نقاش} \fa{از} \fa{گلها} \fa{نقاشی} \fa{میکند}.
\end{quote}

\begin{quote}
\emph{in naqqâsh az golhâ naqqâshi mikonad.}
\end{quote}

\begin{quote}
This painter paints flowers.
\end{quote}

Also a house painter. Plural \emph{naqqâshân} or \emph{naqqâshhâ}.

\begin{cousinweb}
\emph{naqqâsh} comes from Arabic \emph{naqqāš} (EtymDB), the Arabic word for one who paints or engraves, from the same root as \emph{naqsh}.
\end{cousinweb}

\begin{grammarlens}[title={What you've built}]
One more word for telling, explaining and describing at B1.
\end{grammarlens}

\subsection*{Your turn}
\begin{itemize}
\raggedright
  \item \emph{Say it:} \emph{naqqâsh}
  \item \emph{Say it:} \emph{naqqâsh}, once more
  \item \emph{Say it:} a sentence of your own with \emph{naqqâsh}
  \item \emph{Recall:} say the Persian for to play a role, then the Persian for an actor, then say \emph{naqqâsh} again
\end{itemize}

\begin{tcolorbox}[breakable,colback=teal!4,colframe=teal!35!black,title={Before you move on}]
What does \fa{نقاش} mean? (\textquotedblleft{}A painter\textquotedblright{}.) Say it once more.
\end{tcolorbox}
\section[\texorpdfstring{\fa{بوم} (bum)}{بوم (bum)}]{\fa{بوم} (bum) — a canvas (for painting)}

\label{lesson:FA-C525-bum}



\begin{quote}
Before the new one: say the Persian for an actor, then the Persian for a painter.
\end{quote}

\subsection*{You'll want to know: \fa{بوم}}
\textbf{\fa{بوم}} (\emph{bum}) — \textquotedblleft{}a canvas (for painting)\textquotedblright{}.

\begin{quote}
\fa{نقاش} \fa{یک} \fa{بوم} \fa{بزرگ} \fa{خرید}.
\end{quote}

\begin{quote}
\emph{naqqâsh yek bum-e bozorg kharid.}
\end{quote}

\begin{quote}
The painter bought a large canvas.
\end{quote}

Other words are spelled the same: \emph{bum} 'land' in \emph{marz-o-bum} 'homeland', and \emph{bum} 'owl'.

\begin{cousinweb}
EtymDB traces Persian \emph{bum}, 'land, ground', to Middle Persian and on to Old Persian (EtymDB). The painter's \emph{bum}, the ground a picture is laid on, is usually linked with it, but the link is not settled.
\end{cousinweb}

\begin{grammarlens}[title={What you've built}]
One more word for telling, explaining and describing at B1.
\end{grammarlens}

\subsection*{Your turn}
\begin{itemize}
\raggedright
  \item \emph{Say it:} \emph{bum}
  \item \emph{Say it:} \emph{bum}, once more
  \item \emph{Say it:} a sentence of your own with \emph{bum}
  \item \emph{Recall:} say the Persian for an actor, then the Persian for a painter, then say \emph{bum} again
\end{itemize}

\begin{tcolorbox}[breakable,colback=teal!4,colframe=teal!35!black,title={Before you move on}]
What does \fa{بوم} mean? (\textquotedblleft{}A canvas (for painting)\textquotedblright{}.) Say it once more.
\end{tcolorbox}
\section[\texorpdfstring{\fa{قلم} \fa{مو} (qalam mu)}{قلم مو (qalam mu)}]{\fa{قلم} \fa{مو} (qalam mu) — a paintbrush}

\label{lesson:FA-C525-qalam-mu}



\begin{quote}
Before the new one: say the Persian for a painter, then the Persian for a canvas.
\end{quote}

\subsection*{You'll want to know: \fa{قلم} \fa{مو}}
\textbf{\fa{قلم} \fa{مو}} (\emph{qalam mu}) — \textquotedblleft{}a paintbrush\textquotedblright{}.

\begin{quote}
\fa{او} \fa{با} \fa{یک} \fa{قلم} \fa{موی} \fa{نازک} \fa{کار} \fa{میکند}.
\end{quote}

\begin{quote}
\emph{u bâ yek qalam mu-ye nâzok kâr mikonad.}
\end{quote}

\begin{quote}
He works with a thin paintbrush.
\end{quote}

\emph{qalam} 'pen' + \emph{mu} 'hair': a hair pen.

\begin{cousinweb}
\emph{qalam mu} is built in Persian on \emph{qalam}, 'pen', from Arabic \emph{qalam} (EtymDB), and \emph{mu}, 'hair', which EtymDB traces to Middle Persian (EtymDB).
\end{cousinweb}

\begin{grammarlens}[title={What you've built}]
One more word for telling, explaining and describing at B1.
\end{grammarlens}

\subsection*{Your turn}
\begin{itemize}
\raggedright
  \item \emph{Say it:} \emph{qalam mu}
  \item \emph{Say it:} \emph{qalam mu}, once more
  \item \emph{Say it:} a sentence of your own with \emph{qalam mu}
  \item \emph{Recall:} say the Persian for a painter, then the Persian for a canvas, then say \emph{qalam mu} again
\end{itemize}

\begin{tcolorbox}[breakable,colback=teal!4,colframe=teal!35!black,title={Before you move on}]
What does \fa{قلم} \fa{مو} mean? (\textquotedblleft{}A paintbrush\textquotedblright{}.) Say it once more.
\end{tcolorbox}
\section[\texorpdfstring{\fa{آبرنگ} (âbrang)}{آبرنگ (âbrang)}]{\fa{آبرنگ} (âbrang) — watercolour}

\label{lesson:FA-C525-abrang}



\begin{quote}
Before the new one: say the Persian for a canvas, then the Persian for a paintbrush.
\end{quote}

\subsection*{You'll want to know: \fa{آبرنگ}}
\textbf{\fa{آبرنگ}} (\emph{âbrang}) — \textquotedblleft{}watercolour\textquotedblright{}.

\begin{quote}
\fa{دخترم} \fa{با} \fa{آبرنگ} \fa{یک} \fa{دریا} \fa{کشید}.
\end{quote}

\begin{quote}
\emph{dokhtaram bâ âbrang yek daryâ keshid.}
\end{quote}

\begin{quote}
My daughter painted a sea in watercolour.
\end{quote}

\emph{âb} 'water' + \emph{rang} 'colour'.

\begin{cousinweb}
\emph{âbrang} is built on two inherited words: \emph{âb}, 'water', and \emph{rang}, 'colour, dye', both of which EtymDB traces to Middle Persian (EtymDB).
\end{cousinweb}

\begin{grammarlens}[title={What you've built}]
One more word for telling, explaining and describing at B1.
\end{grammarlens}

\subsection*{Your turn}
\begin{itemize}
\raggedright
  \item \emph{Say it:} \emph{âbrang}
  \item \emph{Say it:} \emph{âbrang}, once more
  \item \emph{Say it:} a sentence of your own with \emph{âbrang}
  \item \emph{Recall:} say the Persian for a canvas, then the Persian for a paintbrush, then say \emph{âbrang} again
\end{itemize}

\begin{tcolorbox}[breakable,colback=teal!4,colframe=teal!35!black,title={Before you move on}]
What does \fa{آبرنگ} mean? (\textquotedblleft{}Watercolour\textquotedblright{}.) Say it once more.
\end{tcolorbox}
\section[\texorpdfstring{\fa{گالری} (gâleri)}{گالری (gâleri)}]{\fa{گالری} (gâleri) — a gallery}

\label{lesson:FA-C525-galeri}



\begin{quote}
Before the new one: say the Persian for a paintbrush, then the Persian for watercolour.
\end{quote}

\subsection*{You'll want to know: \fa{گالری}}
\textbf{\fa{گالری}} (\emph{gâleri}) — \textquotedblleft{}a gallery\textquotedblright{}.

\begin{quote}
\fa{نقاشی} \fa{های} \fa{او} \fa{در} \fa{یک} \fa{گالری} \fa{است}.
\end{quote}

\begin{quote}
\emph{naqqâshihâ-ye u dar yek gâleri ast.}
\end{quote}

\begin{quote}
Her paintings are in a gallery.
\end{quote}

\emph{gâleri-ye naqqâshi} — an art gallery.

\begin{cousinweb}
\emph{gâleri} was borrowed from French \emph{galerie} (EtymDB).
\end{cousinweb}

\begin{grammarlens}[title={What you've built}]
One more word for telling, explaining and describing at B1.
\end{grammarlens}

\subsection*{Your turn}
\begin{itemize}
\raggedright
  \item \emph{Say it:} \emph{gâleri}
  \item \emph{Say it:} \emph{gâleri}, once more
  \item \emph{Say it:} a sentence of your own with \emph{gâleri}
  \item \emph{Recall:} say the Persian for a paintbrush, then the Persian for watercolour, then say \emph{gâleri} again
\end{itemize}

\begin{tcolorbox}[breakable,colback=teal!4,colframe=teal!35!black,title={Before you move on}]
What does \fa{گالری} mean? (\textquotedblleft{}A gallery\textquotedblright{}.) Say it once more.
\end{tcolorbox}
